%%=======================================================================
% !Mode:: "TeX:UTF-8"
% !TEX program  = PdfLaTeX
%%=======================================================================
% 模板名称：hitszbeamer
% 模板版本：V1.1.0
% 模板作者：杨敬轩（Jingxuan Yang）
% 联系作者：yangjingxuan@stu.hit.edu.cn & yanglatex2e@gmail.com
% 模板交流：QQ群：1039392552，加群请备注LaTeX、hitszthesis相关说明
% 模板适用：哈尔滨工业大学（深圳）Beamer模板
% 模板编译：手动编译方法参看 README.md 或 hitszbeamer.pdf
%          编译beamer之前必须编译说明文档：make doc 或双击 makedoc.bat
%          编译说明文档同时分离出四个样式文件 *hitszbeamer.sty
%          GNU make 工具：make beamer
%          Windows批处理脚本：双击 makebeamer.bat 自动编译论文
%          更多编译细节详见说明文档：hitszbeamer.pdf
% 更新时间：2022/05/20
% 模板帮助：请**务必务必务必**阅读 hitszbeamer.pdf 说明文档，文档查看方法：
%          下载模板文件夹里就有，如果是从CTAN上安装更新本模板，则通过
%          cmd 命令行：texdoc hitszbeamer
%          推荐前往模板的 GitHub 仓库获取最新文件，地址：
%          https://github.com/YangLaTeX/hitszbeamer
%%=======================================================================

% 设置文档类别为 <beamer>，<compress>选项含义为尽量压缩导航栏
% 添加aspectratio设置成16:9
\documentclass[compress, aspectratio=169]{beamer}
% \documentclass[compress]{beamer}
\PassOptionsToPackage{numbers, square, sort&compress}{natbib} % 解决 Package natbib Error: Bibliography not compatible with author-year citations

% 使用 <hitszbeamer> 主题
\usetheme{hitszbeamer}

% cite命令引用在右上角
\let\citeoriginal\cite
\renewcommand{\cite}[1]{\textsuperscript{\citeoriginal{#1}}}

% 每一章节前插入目录
\AtBeginSection[] 
{
  \begin{frame}
    \frametitle{目录}
    \tableofcontents[currentsection]
  \end{frame}
}

% 开始写文章
\begin{document}


% 图片存放路径
\graphicspath{{figures/}}

% 封面信息
\title[移位异或纠删码的\\高效实现研究]{移位异或纠删码的高效实现研究\\[1mm] \small 开题报告\vspace{-0.7em}}
\author[孙铎]{学生：孙铎 \\[1mm] 导师：王强\ 副教授}
\institute[哈尔滨工业大学（深圳）计算机科学与技术学院]{\small  哈尔滨工业大学（深圳）计算机科学与技术学院}
\date{\small \vskip -10pt 2025年09月12日}

% 标题页
\begin{frame}
		\maketitle
\end{frame}

% 目录页
\section*{目录}
\frame{
  \frametitle{\secname}
  \tableofcontents[sections=all]
}

% 背景动机
\section{研究内容与进度}

\begin{frame}{研究内容与进度}{先看一个真实问题：存储系统会丢块}

    \begin{columns}[T,onlytextwidth]
        \begin{column}{0.48\textwidth}
            \begin{block}{为什么不能只靠三副本}
                \footnotesize
                \begin{itemize}
                    \setlength{\itemsep}{3pt}
                    \item 云存储、文件系统和对象存储都要容忍磁盘或节点失效
                    \item 三副本实现简单，但 1\,TB 数据要占约 3\,TB 空间
                    \item 纠删码把数据切成 $k$ 个数据块，再生成 $m$ 个校验块，用更低冗余恢复丢失块
                \end{itemize}
            \end{block}
        \end{column}

        \begin{column}{0.52\textwidth}
            \begin{block}{本课题关心的具体场景}
                \footnotesize
                \begin{itemize}
                    \setlength{\itemsep}{3pt}
                    \item 存储系统中同时可能丢数据块，也可能丢校验块
                    \item 编码和恢复速度越快，系统写入、修复和降级读越不容易被拖慢
                    \item 目标是在 Ceph 这类真实系统里做出高吞吐纠删码实现
                \end{itemize}
            \end{block}

            \vspace{0.3em}
            \begin{block}{一句话概括}
                \centering
                \footnotesize
                少占空间，同时在块丢失时尽快把数据和校验都恢复回来。
            \end{block}
        \end{column}
    \end{columns}

\end{frame}

\begin{frame}{研究内容与进度}{为什么选择移位异或码}

    \setbeamerfont{caption}{size=\scriptsize}
    \begin{columns}[T,onlytextwidth]
        \column{0.42\textwidth}
        \begin{block}{传统 Reed-Solomon 码}
            \scriptsize
            \begin{itemize}
                \setlength{\itemsep}{1pt}
                \item 可靠、成熟，Ceph 等系统已广泛使用
                \item 编解码依赖有限域矩阵运算
                \item 高性能实现优化复杂
            \end{itemize}
        \end{block}

        \begin{block}{Two-tone 移位异或码}
            \scriptsize
            \begin{itemize}
                \setlength{\itemsep}{1pt}
                \item 校验块由数据块移位后异或得到
                \item 移位可用地址偏移实现
                \item 理论上很快，但缺少工程级高性能实现
            \end{itemize}
        \end{block}

        \column{0.58\textwidth}
        \centering
        \scriptsize\textbf{例：4 个数据块经过不同偏移后异或生成 3 个校验块}

        \vspace{0.2em}
        \includegraphics[width=0.96\textwidth,height=0.49\textheight,keepaspectratio]{Example of Two-tone Encoding.pdf}
    \end{columns}

\end{frame}

\begin{frame}{研究内容与进度}{中期已经做出了什么}

    \begin{columns}[T,onlytextwidth]
        \column{0.50\textwidth}
        \begin{block}{已经完成的系统}
            \footnotesize
            \begin{itemize}
                \setlength{\itemsep}{3pt}
                \item 完成 Two-tone 编码、解码和修复流程分析
                \item 实现约 2000 行 C++ FastTT 核心库
                \item 把 FastTT 做成 Ceph v18.2.7 纠删码插件
                \item 通过 Ceph 接口测试和集群端到端读写验证
            \end{itemize}
        \end{block}

        \column{0.50\textwidth}
        \begin{block}{已经得到的结果}
            \footnotesize
            \begin{itemize}
                \setlength{\itemsep}{3pt}
                \item 相较 ISA-L：编码平均提升 \textbf{59.1\%}
                \item 相较 ISA-L：解码平均提升 \textbf{50.9\%}
                \item 单块丢失恢复最高提升 \textbf{182.1\%}
                \item 与学术基线 Cerasure 对比也取得整体优势
                \item 相关成果已整理为论文并投稿 IPDPS 2026
            \end{itemize}
        \end{block}
    \end{columns}

\end{frame}


% 文献综述
\section{国内外研究现状及分析}

\begin{frame}{国内外研究现状及分析}
    \textbf{RS码理论研究综述：早期奠定基础的编解码理论}
    
    \vspace{0.5em}

    \begin{enumerate}
        \item 由Reed和Solomon提出\cite{reedPolynomialCodesCertain1960}，基于有限域Vandermonde矩阵乘法。
        \item Luby等人将乘法变换为异或，并使用求逆更快的Cauchy矩阵\cite{lubyXORBasedErasureResilientCoding1995}。
    \end{enumerate}
    
    这种理论最优的RS码是后续优化研究的基础。

    \vspace{0.5em}

    \textbf{RS码优化研究综述：重点研究工作及其优化方案总结}
    
    \vspace{0.5em}

    \begin{table}[]
    \centering
    \begin{tabular}{|c|c|}
    \hline
    减少矩阵中$1$的个数进而减少运算 & Zerasure\cite{zhouFastErasureCoding2020}、Cerasure\cite{wangFastAccelerationStrategies2025} \\ \hline
    提取并复用公共异或运算表达式    & SLPEC\cite{uezatoAcceleratingXORbasedErasure2021}、Cerasure                                 \\ \hline
    内存与缓存访问优化  & Zerasure、SLPEC、Cerasure \\ \hline
    降低数据指针相关开销 & Cerasure                \\ \hline
    \end{tabular}
    \end{table}

    Cerasure综合各优化策略，并提出指针开销优化，是RS码SOTA实现。

\end{frame}

\begin{frame}{国内外研究现状及分析}

    \textbf{移位异或纠删码理论研究综述：从ZigZag到Two-tone}
    
    \vspace{1em}

    \begin{columns}[t]
        \column{0.5\textwidth}
        ZigZag Decodable Codes\cite{gongZigzagDecodableCodes2017}是较早提出的移位异或纠删码，通过构建特殊的移位矩阵实现了仅依赖异或和移位的高效编解码
        \column{0.5\textwidth}
        twotone
    \end{columns}
    
\end{frame}


% 研究方案
\section{已完成的主要研究工作}

\subsection{原理瓶颈与优化框架}

\begin{frame}{已完成的主要研究工作}{Two-tone编解码原理与性能瓶颈}

    \setbeamerfont{caption}{size=\scriptsize}
    \begin{columns}[T]
        \column{0.46\textwidth}
        \footnotesize
        \begin{block}{Two-tone 的基本特征}
            \footnotesize
            \begin{itemize}
                \setlength{\itemsep}{2pt}
                \item 编码仅由移位和异或构成，移位由地址偏移实现
                \item 解码可抽象为代入、消除和修复过程
                \item 核心瓶颈是\textbf{内存访问}，不是 XOR 算术
            \end{itemize}
        \end{block}

        \column{0.54\textwidth}
        \begin{figure}
            \centering
            \includegraphics[width=\textwidth,height=0.50\textheight,keepaspectratio]{Example of Two-tone Decoding.pdf}
            \vspace{-0.7em}
            \caption{Two-tone解码举例}
        \end{figure}
        \vspace{-0.4em}
        \footnotesize
        \begin{itemize}
            \setlength{\itemsep}{2pt}
            \item 冗余访存和时间局部性差
            \item 横向/纵向遍历难以同时适配所有场景
            \item 解码与修复分离导致额外遍历
        \end{itemize}
    \end{columns}

\end{frame}

\begin{frame}[shrink=10]{已完成的主要研究工作}{FastTT整体优化框架}

    \setbeamerfont{caption}{size=\scriptsize}
    \begin{columns}[T]
        \column{0.40\textwidth}
        \begin{block}{设计原则}
            \small 不改变 Two-tone 数学结构，围绕\textbf{内存访问路径}做系统级优化。
        \end{block}
        \vspace{0.2em}
        \footnotesize
        \begin{block}{四项协同优化}
            \footnotesize
            \begin{itemize}
                \setlength{\itemsep}{3pt}
                \item \textbf{EPM}：统一数据/校验恢复
                \item \textbf{CFWP}：窗口化提升缓存局部性
                \item \textbf{NMA}：相邻块合并降低写放大
                \item \textbf{TPS}：按丢失模式选择遍历路径
            \end{itemize}
        \end{block}
        \centering
        \footnotesize 编码侧：CFWP+NMA\quad 解码侧：四项协同

        \column{0.60\textwidth}
        \centering
        \includegraphics[width=0.88\textwidth,height=0.42\textheight,keepaspectratio]{Two-tone Optimization Overview.pdf}
        \vspace{0.2em}
        \scriptsize FastTT整体优化框架
    \end{columns}

\end{frame}

\subsection{四项优化}

\begin{frame}{已完成的主要研究工作}{EPM 与 CFWP：减少冗余遍历并提升局部性}

    \footnotesize
    \begin{columns}[T]
        \column{0.5\textwidth}
        \begin{block}{EPM：统一解码抽象}
            \footnotesize
            \begin{itemize}
                \setlength{\itemsep}{2pt}
                \item 将丢失数据块映射到存活校验块，丢失校验块映射到自身
                \item 把校验块修复融入解码流程，省去独立 repair 遍历
            \end{itemize}
        \end{block}
        \centering
        \includegraphics[width=0.84\textwidth,height=0.34\textheight,keepaspectratio]{Erasure-Parity Mapping Decoding Demonstration.pdf}
        \vspace{0.2em}
        \scriptsize 作用：减少冗余访存，统一处理所有丢失组合

        \column{0.5\textwidth}
        \begin{block}{CFWP：缓存友好窗口划分}
            \footnotesize
            \begin{itemize}
                \setlength{\itemsep}{2pt}
                \item 将块切分为窗口，窗口内完成相关移位异或
                \item 首次访问时初始化，避免全量预清零带来的额外开销
            \end{itemize}
        \end{block}
        \centering
        \includegraphics[width=0.92\textwidth,height=0.34\textheight,keepaspectratio]{Cache-Friendly Window Demonstration.pdf}
        \vspace{0.2em}
        \scriptsize 作用：在数据离开缓存前完成更多有效计算
    \end{columns}

\end{frame}

\begin{frame}{已完成的主要研究工作}{NMA 与 TPS：降低写放大并加速常见故障}

    \footnotesize
    \begin{columns}[T]
        \column{0.5\textwidth}
        \begin{block}{NMA：相邻合并访问}
            \footnotesize
            \begin{itemize}
                \setlength{\itemsep}{2pt}
                \item 相邻数据块先在 SIMD 寄存器内异或合并
                \item 合并后再写回校验块，减少校验块写入次数
            \end{itemize}
        \end{block}
        \centering
        \includegraphics[width=0.94\textwidth,height=0.36\textheight,keepaspectratio]{Neighbor-Merging Access Pattern.pdf}
        \vspace{0.2em}
        \scriptsize 作用：降低写放大，提升编码端吞吐

        \column{0.5\textwidth}
        \begin{block}{TPS：遍历模式选择}
            \footnotesize
            \begin{itemize}
                \setlength{\itemsep}{2pt}
                \item 单块丢失直接走纵向单遍快路径
                \item 避免完整代入/消除流程，每码字仅访问一次
            \end{itemize}
        \end{block}
        \centering
        \includegraphics[width=0.84\textwidth,height=0.36\textheight,keepaspectratio]{Vertical Traversal Pattern Demonstration.pdf}
        \vspace{0.2em}
        \scriptsize 作用：显著加速最常见的单块失效场景
    \end{columns}

\end{frame}


% 目标进度
\section{实验评估与结果}

\begin{frame}{实验评估与结果}{编码与解码吞吐}

    \footnotesize
    环境：i9-14900K + AVX2、64\,GB DDR5、块2\,MB；基线：Jerasure、\textbf{ISA-L}（Ceph生产级）、\textbf{Cerasure}（学术SOTA）。

    \vspace{0.4em}
    \begin{columns}[T]
        \column{0.5\textwidth}
        \begin{block}{编码吞吐（GB/s）}
            \centering \scriptsize
            \begin{tabular}{cccc}
                \toprule
                配置 & FastTT & ISA-L & Jera. \\
                \midrule
                (6, 2)  & \textbf{34.74} & 24.37 & 17.03 \\
                (8, 3)  & \textbf{31.35} & 17.03 & 7.08  \\
                (10, 4) & \textbf{17.59} & 13.56 & 4.58  \\
                \bottomrule
            \end{tabular}
        \end{block}
        \centering 相较ISA-L平均\textbf{+59.1\%}，Jerasure \textbf{+261\%}

        \column{0.5\textwidth}
        \begin{block}{解码吞吐（GB/s）}
            \centering \scriptsize
            \begin{tabular}{ccccc}
                \toprule
                配置 & 丢失 & FastTT & ISA-L & Jera. \\
                \midrule
                (6, 2)  & 1 & \textbf{95.99} & 38.38 & 42.70 \\
                (8, 3)  & 1 & \textbf{93.11} & 33.01 & 37.04 \\
                (8, 3)  & 3 & \textbf{21.44} & 18.38 & 5.00  \\
                \bottomrule
            \end{tabular}
        \end{block}
        \centering 相较ISA-L平均\textbf{+50.9\%}，单块丢失最高\textbf{+182\%}
    \end{columns}

    \vspace{0.4em}
    \centering 对比学术SOTA Cerasure：编码\textbf{+13.2\%}、解码\textbf{+16.0\%}（最高+43.6\%）。

\end{frame}

\begin{frame}{实验评估与结果}{消融实验与端到端吞吐}

    \footnotesize
    \begin{columns}[T]
        \column{0.55\textwidth}
        \begin{block}{消融实验 $(8,3)$（GB/s）}
            \centering \scriptsize
            \begin{tabular}{ccccrr}
                \toprule
                CFWP & NMA & EPM & TPS & 编码 & 解码 \\
                \midrule
                           &            &            &            & 13.16 & 7.54  \\
                           &            & \checkmark &            & 13.16 & 16.85 \\
                \checkmark &            & \checkmark &            & 23.32 & 24.10 \\
                \checkmark & \checkmark & \checkmark &            & 31.23 & 24.10 \\
                \checkmark & \checkmark & \checkmark & \checkmark & \textbf{32.78} & \textbf{35.23} \\
                \bottomrule
            \end{tabular}
        \end{block}
        \centering \scriptsize 四项优化针对不同瓶颈，贡献清晰且\textbf{可叠加}

        \column{0.45\textwidth}
        \begin{block}{端到端 \emph{PUT} 吞吐（MB/s）}
            \centering \scriptsize
            \begin{tabular}{cccc}
                \toprule
                配置 & FastTT & ISA-L & Jera. \\
                \midrule
                (6, 2)  & \textbf{457.7} & 444.6 & 401.4 \\
                (8, 3)  & \textbf{877.3} & 817.6 & 709.3 \\
                (10, 4) & \textbf{573.0} & 546.4 & 490.1 \\
                \bottomrule
            \end{tabular}
        \end{block}
        \centering \scriptsize 含磁盘I/O仍\textbf{持续优于ISA-L}（平均+7.1\%）
    \end{columns}

    \vspace{0.5em}
    \centering
    \textbf{编/解码、消融与端到端实验全面验证优化有效性，达成开题核心目标。}

\end{frame}


% \section{参考文献}
\bibliographystyle{hitszbeamer}
\begin{frame}{参考文献}
  \tiny
  \setlength{\bibsep}{-0.5pt}
  \bibliography{reference}
\end{frame}

% Q&A
% \section{Q\&A}
\begin{frame}{\secname~ }
	\begin{center}
 		\huge {感谢各位老师的耐心观看\\请老师批评指正}\\
		\vspace{1cm}
		\huge {Q\&A}
	\end{center}
\end{frame}



% 结束文档撰写
\end{document}

